\section{Introduction and source comparison}\label{sec:intro}
With zero initial data the zero control makes the terminal temperature zero. The point considered here is instead the existence of many \emph{nonzero} controls with the same zero temperature. This is a non-injectivity question for an output map, not null controllability from arbitrary initial states and not full-state return to rest. The controls considered below belong to the stated smooth compatible class.

Fattorini and Russell's 1971 paper \cite{FattoriniRussell1971} uses exponential moments for a one-dimensional scalar parabolic equation; its Theorems 3.3--3.4 treat reachable and null-controllable states including smooth compactly supported controls under the stated decay hypotheses. Their 1974 Theorems 1.2 and 1.5(b) \cite{FattoriniRussell1974} give the fixed-time negative-exponential biorthogonal estimate used here. Neither source directly addresses the coupled slow branch. The 2011 theorem of Tenenbaum and Tucsnak \cite{TenenbaumTucsnak2011} gives arbitrary-initial-data null controllability and cost estimates under a self-adjoint spectral hypothesis and a spectral inequality. Its eigenvalues tend to infinity; this does not cover our full accumulating spectrum directly.

Ivanov and Pandolfi \cite{IvanovPandolfi2009} study rest from arbitrary initial states for a heat equation in which the memory acts on the spatial Laplacian. Their Obstruction 1 uses accumulation of characteristic roots and analyticity to obstruct interpolation. We credit that mechanism when establishing that zero temperature and zero memory together force zero control from zero start; the conclusion is not claimed as new. Our bounded ODE coupling yields a different equation and target, and their negative rest theorem is not invoked as a theorem about our model. Their paper also gives conditional negative results through a completeness obstruction and fractional-kernel arguments.

Guerrero and Imanuvilov \cite{Guerrero2013} show non-null-controllability for some nonzero initial data in their specified heat--memory equations, including a model with memory acting on the state itself, using adjoint observability. Halanay and Pandolfi \cite{HalanayPandolfi2015}, for the Coleman--Gurtin equation and their stated control and terminal-trace hypotheses, obtain approximate temperature controllability and failure of temperature null controllability for some initial data when the relevant resolvent is nonzero. Neither statement settles the exact zero-start memory image below. Ammar Khodja et al.~\cite{Ammar2014} relate minimal time to condensation for a Riesz-basis spectral problem under sector and reciprocal-summability assumptions; the bounded accumulating branch here violates the latter for the combined spectrum. Shklyar's preprint \cite{Shklyar2024} studies partial null controllability by moments under a no-finite-accumulation spectrum assumption. Neither result directly settles our problem. These targeted comparisons do not prove originality or exhaust the literature.

The following sections formulate the frozen model, establish the paired moments, and give an abstract kernel construction. The memory trace restrictions are necessary conditions, not an exact characterization or a density theorem. The construction and its physical interpretation remain subject to independent review.
